% Generated by scripts/generate_protocol_appendix.py.
% Source: experiments/protocols/antidote-feasibility-v1.1.json.

\noindent\textbf{Protocol identity.} \texttt{ANT-PROT-FEAS-001}, version~1.1.0, was frozen as a \texttt{frozen-design-protocol} artifact on 2026-09-07. Its byte-exact SHA-256 is shown below in two contiguous halves.
\begin{center}
  \small\texttt{8d6848f148a627424c566f0438110377}\\[-0.2em]
  \small\texttt{9debd65b4563bd4712652e09cd715024}
\end{center}

\noindent\textbf{Protocol succession.} This is the pre-collection corrective successor to version~1.0.0. The superseded protocol remains preserved at \nolinkurl{experiments/protocols/antidote-feasibility-v1.json} with byte-exact SHA-256 \nolinkurl{dbedbcbb74303373a7e00e95ad063b025d8cf6033864f4ad5a390d15703ba403}. No qualifying record or human collection began under the predecessor.

\noindent\textbf{Authority status.} Collection authority is \textbf{false}; no formal human study has begun. The amendment and deviation log contains 0 entries.

\subsection*{Evidence stages}
\begin{description}[leftmargin=3.3em,style=nextline,itemsep=0.45em]
  \item[\texttt{D0}] \textbf{Design traceability.}\\ Current state: \nolinkurl{available-for-design-review}. Unit: one governed transformation or authority boundary.
  \item[\texttt{T0}] \textbf{Synthetic conformance and recovery.}\\ Current state: \nolinkurl{executable-mock-slice}. Unit: one deterministic fixture, command trace, or injected failure path.
  \item[\texttt{T1}] \textbf{Real-model technical qualification.}\\ Current state: \nolinkurl{blocked-no-real-model}. Unit: one immutable generation specification and terminal generation result.
  \item[\texttt{H1}] \textbf{Repeated-session within-person feasibility.}\\ Current state: \nolinkurl{blocked-no-collection-authority}. Unit: one assigned session attempt, with exposure analyzed only when verified audio was actually played.
\end{description}

\subsection*{Prospective H1 schedule}
\begin{itemize}[leftmargin=1.35em,itemsep=0.4em]
  \item 6 blocks $\times$ 3 conditions = 18 planned completed exposures, with an absolute ceiling of 24 scheduled attempts.
  \item Each proposed exposure is 8 minutes; at least 48 hours separate exposures, with at most 3 in seven days.
  \item The later aftereffect window targets 24 hours and accepts 18--30 hours.
  \item The balanced block-order assignment contains exactly 720 schedules ($6!$), selected by the frozen uniform algorithm; the six H1 block seeds are committed before session~1.
\end{itemize}

\begin{description}[leftmargin=2.8em,style=nextline,itemsep=0.45em]
  \item[\texttt{G}] \textbf{generic generation.} Create an eight-minute instrumental ambient listening piece with no lyrics, sudden transients, or intentionally intense effects.
  \item[\texttt{S}] \textbf{structured non-personalized journey.} Use the selected current-state and desired-transition categories with one frozen meet-hold-transition-integrate template; do not use free-text history, personal descriptors, or prior-session mappings.
  \item[\texttt{P}] \textbf{negotiated personal semantic journey.} Use only the explicitly approved working projection, person-authored semantic mixins, inclusions, exclusions, and accepted journey-plan revision.
\end{description}

\subsection*{Human-collection activation gates}
\begin{itemize}[leftmargin=1.35em,itemsep=0.42em]
  \item \texttt{GATE-REAL-MODEL} --- \textbf{not satisfied.} Real model, adapter, artifact, cancellation, and T1 technical qualification complete.
  \item \texttt{GATE-PRIVACY} --- \textbf{not satisfied.} Encryption, retention, deletion, backup, projection invalidation, minimum-data private research export, and separately minimized public audit-envelope controls reviewed for the intended setting.
  \item \texttt{GATE-CONSENT} --- \textbf{not satisfied.} The version 2 scope-specific consent contract and participant-facing consent, withdrawal, adverse-response, and data-use materials are implemented and approved for the actual study context.
  \item \texttt{GATE-INDEPENDENT-REVIEW} --- \textbf{not satisfied.} An independent qualified reviewer records the applicable ethics, safety, and researcher-participant determination; no approval is presumed.
  \item \texttt{GATE-ASSIGNMENT} --- \textbf{not satisfied.} Assignment seed, all 720 eligible schedules, selected complete schedule, and the six fixed block seeds are frozen and their enforcement is verified before session 1.
  \item \texttt{GATE-INSTRUMENT} --- \textbf{not satisfied.} The timed session workflow, exact prompts and anchors, version 2 response and append-only correction contract, playback setup, accessibility path, and implementation match this protocol.
  \item \texttt{GATE-ANALYSIS} --- \textbf{not satisfied.} Analysis code and synthetic validation fixtures reproduce the frozen estimands, denominators, missingness, and sensitivity rules.
  \item \texttt{GATE-STOP-GO} --- \textbf{not satisfied.} A documented pilot stop/go decision authorizes only the explicitly reviewed next step.
\end{itemize}

\subsection*{Freeze and audit rule}
\begin{itemize}[leftmargin=1.35em,itemsep=0.42em]
  \item Any material change creates a new semantic protocol version, a new content lock, an updated analysis plan, and renewed review before collection.
  \item Never rewrite the frozen file. Append a timestamped amendment or deviation with rationale, affected records, evidence impact, author, reviewer, and disposition.
  \item Report zero movement, null contrasts, opposite-direction movement, mismatch, unwanted intensity, burden, harm, stops, failures, and adverse events with the same visibility as favorable observations.
\end{itemize}
